\subsection{\texorpdfstring{The maps \(f, v\) and \(\Phi\)}{The maps f, v, and Φ}}

Let A be a ring. Endow \(\mathbf{A}^{\mathbf{N}}\) with the structure of the product ring. Denote by \(f_{\mathbf{A}}\), or simply \(f\), the endomorphism \((a_{n})_{n\in\mathbf{N}}\mapsto(a_{n+1})_{n\in\mathbf{N}}\) of \(\mathbf{A}^{\mathbf{N}}\). Denote by \(v_{\mathbf{A}}\), or simply \(v\), the endomorphism of the additive group underlying \(\mathbf{A}^{\mathbf{N}}\) which sends \((a_{n})_{n\in\mathbb{N}}\) to \((0,p\cdot a_{0},p\cdot a_{1},...)\).

For every integer \(m \geqslant 0\), denote by \(\Phi_{m}\) the map from \(\mathbf{A}^{\mathbf{N}}\) into A which sends \(\boldsymbol{a} = (a_{n})_{n \in \mathbf{N}}\) to \(\Phi_{m}(a_{0}, \ldots, a_{m})\). We denote by \(\Phi_{\mathrm{A}}\), or simply \(\Phi\), the map \(\boldsymbol{a} \mapsto (\Phi_{n}(\boldsymbol{a}))_{n \in \mathbf{N}}\) of \(\mathbf{A}^{\mathbf{N}}\) into itself.

Lemma 2.--- Let A be a ring endowed with an endomorphism σ satisfying σ(a) ≡ a\^{}p mod. p.A for every a ∈ A. Let n ≥ 1 be an integer and let a\_0, \ldots, a\_\{n-1\} be elements of A. Set u\_i = Φ\_i(a\_0, \ldots, a\_i) for 0 ≤ i ≤ n - 1. Let u\_n be an element of A. The following conditions are equivalent :

\begin{enumerate}
\def\labelenumi{\alph{enumi})}
\item
  There exists \(a_{n} \in A\) such that \(u_{n} = \Phi_{n}(a_{0}, \ldots, a_{n})\).
\item
  \(\sigma(u_{n-1}) \equiv u_n \pmod{p^nA}\).
\end{enumerate}

For \(0 \leqslant i \leqslant n - 1\), one has \(\sigma(a_i) \equiv a_i^p \pmod{pA}\). By Proposition 1 of no. 1, applied to the case where \(\mathbf{J}_k = p^kA\) for \(k \in \mathbf{N}\) and where \(m = 1\), we have the congruence

\[
\Phi_ {n - 1} \left(\sigma \left(a _ {0}\right), \dots , \sigma \left(a _ {n - 1}\right)\right) \equiv \Phi_ {n - 1} \left(a _ {0} ^ {p}, \dots , a _ {n - 1} ^ {p}\right) \pmod{p ^ {n}A},\tag{7}
\]

that is,

\[
\sigma (u _ {n - 1}) \equiv \Phi_ {n - 1} (a _ {0} ^ {p},..., a _ {n - 1} ^ {p}) \pmod{p ^ {n}A}.\tag{8}
\]

Now, according to formula (2), the relation \(u_n = \Phi_n(a_0, \ldots, a_n)\) is equivalent to

\[
u _ {n} = \Phi_ {n - 1} (a _ {0} ^ {p},..., a _ {n - 1} ^ {p}) + p ^ {n} a _ {n}.\tag{9}
\]

The lemma follows from this.

PROPOSITION 2. --- Let A be a ring.

\begin{enumerate}
\def\labelenumi{\alph{enumi})}
\item
  If \(p.1_{\mathbf{A}}\) is not a zero divisor in \(\mathbf{A}\) , the map \(\Phi_{\mathbf{A}}\) is injective.
\item
  If \(p.1_{\mathbf{A}}\) is invertible in \(\mathbf{A}\) , the map \(\Phi_{\mathbf{A}}\) is bijective.
\item
  If \(\sigma\) is an endomorphism of the ring A, satisfying \(\sigma(a) \equiv a^p \bmod pA\) for all \(a \in \mathbf{A}\), the image \(\mathbf{A}'\) of \(\Phi_{\mathbf{A}}\) is a subring of \(\mathbf{A}^{\mathbf{N}}\), stable under \(f_{\mathbf{A}}\) and \(v_{\mathbf{A}}\). It is the set of elements \((u_n)_{n \in \mathbf{N}}\) of \(\mathbf{A}^{\mathbf{N}}\) such that \(\sigma(u_n) \equiv u_{n+1} \bmod p^{n+1}A\) for all \(n \in \mathbf{N}\).
\end{enumerate}

If \(\boldsymbol{a} = (a_n)_{n \in \mathbb{N}}\) and \(\boldsymbol{u} = (u_n)_{n \in \mathbb{N}}\) are elements of \(A^{\mathbb{N}}\) , the relation \(\Phi_A(\boldsymbol{a}) = \boldsymbol{u}\) is equivalent, according to formula (2), to the equalities

\[
\left\{ \begin{array}{l} u _ {0} = a _ {0}, \\ u _ {n} = \Phi_ {n - 1} \left(a _ {0} ^ {p}, \dots , a _ {n - 1} ^ {p}\right) + p ^ {n} a _ {n} \quad \text {for all} \quad n \geqslant 1. \end{array} \right.\tag{10}
\]

Let \(\boldsymbol{u} = (u_{n})_{n \in \mathbb{N}}\) be an element of \(A^{\mathbb{N}}\). When \(p1_{A}\) is a non-zero-divisor in A (resp. when \(p1_{A}\) is invertible in A), there exists at most one sequence \((a_{n})_{n \in \mathbb{N}}\) in A (resp. exactly one sequence \((a_{n})_{n \in \mathbb{N}}\) in A) satisfying the equalities (10), whence a) and b).

We prove \(c)\). By Lemma 2, the image \(\mathbf{A}'\) of \(\Phi_{\mathbf{A}}\) is the set of \(\boldsymbol{u} = (u_n)_{n \in \mathbb{N}}\) in \(\mathbf{A}^{\mathbf{N}}\) such that \(\sigma(u_n) \equiv u_{n+1} \bmod p^{n+1}A\) for all \(n \in \mathbb{N}\). It follows immediately that \(\mathbf{A}'\) is a subring of \(\mathbf{A}^{\mathbf{N}}\), stable under \(f_{\mathbf{A}}\) and \(v_{\mathbf{A}}\).

Remark. — Let \(\boldsymbol{a} = (a_n)_{n \in \mathbb{N}}\) and \(\boldsymbol{u} = (u_n)_{n \in \mathbb{N}}\) be elements of \(A^{\mathbf{N}}\) such that \(\boldsymbol{u} = \Phi_{\mathbf{A}}(\boldsymbol{a})\), and let \(m\) be an integer \(\geqslant 0\). From (10), we deduce the following assertions:

If the \(u_n\), for \(0 \leqslant n \leqslant m\), belong to a subring B of A and if, for every \(x \in A\), the relation \(px \in B\) implies \(x \in B\), then the \(a_n\), for \(0 \leqslant n \leqslant m\), belong to B.

If A is equipped with an \(\mathbf{N}\)-grading, if \(p1_{A}\) is a non-zero-divisor in A, if \(d \in \mathbf{N}\), and if \(u_{n}\) is homogeneous of degree \(dp^{n}\) for \(0 \leqslant n \leqslant m\), then \(a_{n}\) is homogeneous of degree \(dp^{n}\) for \(0 \leqslant n \leqslant m\).

